\documentclass{article}
\usepackage{graphicx} % Required for inserting images

\title{Bachelor Thesis}
\author{Benedikt Adler}
\date{August 2026}

\begin{document}

\maketitle

\section{Introduction}

\section{Forschungsfrage}
Can BunDLe-Net identify aging-caused delays and dedifferentiations in
high-level visual representations

\section{Inhalt}
The project consists of the following initial milestones: \\

    * Initial Analysis of the dataset, data preprocessing (i.e.
    filtering, artifact elimination if necessary) and loading of EEG data
    in the BunDLe-Net model \\

    * Training of the model and evaluation of the results: is the
    hypothesis correct or not? \\

Note that this project is explorative, which means that even though we
think that BunDLe-Net could be sensible in this setting, we have no
proof of that.

\section{Aufbau}
\noindent{Bis Ende August: 
Anschauen von den Daten, heruasfiltern was wir brauchen, überlegen wie man sie auf BunDLe-Net anwenden kann\\
\\
Bis Ende September:
Grundsätzliche Pipeline überlegen und skizzieren\\
\\
Bis Ende Oktober: 
Pipeline implementieren und dokumentieren, eventuelle Verbesserungen finden und testen\\
\\
Bis Ende November:
Pipeline fix fertig in BunDLe-Net integrieren und ausführen, ergebnisse dokumentieren und analysieren\\
\\
Bis Ende Dezember: 
Gesamtbild anschauen und auswerten, Rücksprache mit Luca, schriftlichen Teil vervollständigen\\
\\
Bis Ende Jänner:
Arbeit fertig machen, Code clean schreiben, schriftlichen Teil korrektur lesen lassen und verbessern\\
\\
Februar: 
Abgeben und beten
}

\section{Pipeline}
\begin{enumerate}
    \item EEG in ein Array laden (T, 64), T: Anzahl der Zeitpunkte (also die Zeit die vergangen ist in Sekunden oder ms), N: Anzahl der Kanäle/Neuronen (in unserem Fall 64)
    \item Die Kanäle für alle Personen in derselben Reihenfolge herausfiltern
    \item Die Eventmarke auf eine Label-Zeitreihe B abbilden
    \item Die Eventklassen auf 0 bis 7 reduzieren oder die letzte Dense-Schicht anpassen
    \item X und B zeitlich exakt synchronisieren (also Neuronen aktivität mit Label)
    \item Danach prep\_data(X, B, win=15) aufrufen
\end{enumerate}

Bundle Net Continous anschauen, ob es besser geeignet ist für EEG Daten \\


Für BunDLe-Net gilt daher:

\begin{itemize}
    \item X = EEG-Signal über die Zeit, also Kanäle × Samples.
    \item Die Events sind zunächst nur zeitlich punktuelle Marker.
    \item B kann daraus abgeleitet werden, wenn du EEG-Epochen klassifizieren möchtest.
    \item Die Events sind weder 'Inaktivität' noch direkt kontinuierliche Zuständ
\end{itemize}

Eine sinnvolle Gruppierung für state\_names wäre beispielsweise:

\begin{itemize}
    \item Fixation / kein Stimulus
    \item Gesicht
    \item Tier
    \item Ort
    \item Objekt
    \item Catch-Trial
    \item Button-Press
\end{itemize}

\section{Problems}
The architecture of BunDLe-Net is way to badly optimized for this kind of data load. For just one subject the Computer would need storage of 60 GiB to store one array. 
In order to fix this, I will reduce from float64 to float32. Also I will need to introduce a window mechanism, in order to sequentially send one part of the array to model.tau. 
\end{document}
